\section{Experiments}

In our experiments, we (1) analyze the performance of each individual component of \sysname, (2) evaluate full task performance in both the ``Oracle abstract'' and ``Open'' settings, (3) present promising results verifying claims about COVID-19 using \sysname, and (4) discuss some modeling challenges presented by the dataset.

% TODO add this if time permits.
% See Appendix \todo for versions of all result tables with uncertainty estimates included. These estimates do not affect the findings reported in the main body.

% \paragraph{\sysname\ Components and Details}
% For the \componentone module, \sysname retrieves the top $k=3$ documents as described in \S\ref{sec:baselines}. We use $k=3$ because $k > 3$ reduces end-to-end performance by introducing too many ``false positive'' candidate abstracts. %The recall for evidence abstracts is 68.5\dw{Maybe put in table?}.

% For \componenttwo blah blah

% For \componentthree blah blah.

\subsection{Pipeline components} \label{sec:pipeline_components}
We examine the effects of different training datasets, sentence encoders, and model inputs on the performance of the \componenttwo and \componentthree modules. See Appendix \ref{sec:model_details} for implementation details.


\begin{table}[t]
  \footnotesize
  \setlength{\tabcolsep}{0.5em}
  \centering

  \begin{tabularx}{\columnwidth}{l *{3}{c}  *{3}{c}}
    \toprule
 & \multicolumn{3}{c}{\textsc{Rational-Select.}} & \multicolumn{3}{c}{\textsc{Label-Pred.}} \\
%  & \multicolumn{3}{c}{\textsc{Selection}}  & \multicolumn{3}{c}{\textsc{Prediction}} \\
    % \cmidrule(lr){2-4} \cmidrule(lr){5-7}

    \midrule
    \textbf{Training data} & P & R & F1 & \multicolumn{3}{c}{\textsc{Acc.}} \\

    \arrayrulecolor{black!30}\midrule
    % \hdashline[0.4pt/2pt]

    \fever         &      41.5 &   57.9 &  48.4 &  \multicolumn{3}{c}{67.6} \\
    \snopes        &      42.5 &   62.3 &  50.5 &  \multicolumn{3}{c}{71.3} \\
    \ours          &      73.7 &   70.5 &  \textbf{72.1} &  \multicolumn{3}{c}{75.7} \\
    \fever + \ours &      72.4 &   67.2 &  69.7 &  \multicolumn{3}{c}{\textbf{81.9}} \\

    \arrayrulecolor{black}
    \midrule

    \textbf{Sentence encoder} & P & R & F1 & \multicolumn{3}{c}{\textsc{Acc.}} \\
    \arrayrulecolor{black!30}\midrule
    % \hdashline[0.4pt/2pt]

    \scibert      &      74.5 &   74.3 &  \textbf{74.4} &  \multicolumn{3}{c}{69.2} \\
    \bioroberta   &      75.3 &   69.9 &  72.5 &  \multicolumn{3}{c}{71.7} \\
    \robertabase  &      76.1 &   66.1 &  70.8 &  \multicolumn{3}{c}{62.9} \\
    \robertalarge &      73.7 &   70.5 &  72.1 &  \multicolumn{3}{c}{\textbf{75.7}} \\

    \arrayrulecolor{black}
    \midrule

    \textbf{Model inputs} & P & R & F1 & \multicolumn{3}{c}{\textsc{Acc.}} \\
    \arrayrulecolor{black!30}\midrule
    % \hdashline[0.4pt/2pt]
    Claim-only         &      - &   - &  - &  \multicolumn{3}{c}{44.5} \\
    Abstract-only         &   60.1 &   60.9 &  60.5 &  \multicolumn{3}{c}{53.3} \\

    \arrayrulecolor{black}
    \bottomrule
  \end{tabularx}

  \caption{Comparison of different training datasets, encoders, and model inputs for \componenttwo and \componentthree, evaluated on the \ours dev set. A claim-only model cannot select rationales.}
  \label{tbl:components}

\end{table}
\vspace{.2cm} \noindent{\bf Training Data}
We train on (1) \fever, (2) \snopes, (3) \ours, and (4) \fever pretraining followed by \ours fine-tuning. \robertalarge \cite{Liu2019RoBERTaAR} is used as the sentence encoder.

\vspace{.2cm} \noindent{\bf Sentence encoder}
We fine-tune \scibert \cite{Beltagy2019SciBERTPC}, \bioroberta \cite{Gururangan2020DONTST}, \robertabase, and \robertalarge. \ours is used as training data.

\vspace{.2cm} \noindent{\bf Model Inputs}
We examine the performance of ``claim-only'' and ``abstract-only'' models trained on \ours, using \robertalarge as the sentence encoder. The claim-only model makes label predictions based on the claim text alone, without access to evidence abstracts. The abstract-only model selects rationale sentences and makes label predictions without access to the claim.

\begin{table*}[t]
  \footnotesize
  \setlength{\tabcolsep}{0.5em}
  \centering

  \begin{tabularx}{\linewidth}{*{3}{l} *{6}{c} *{1}{c} *{6}{c}}

    \toprule

    & & & \multicolumn{6}{c}{\textbf{Sentence-level}} & & \multicolumn{6}{c}{\textbf{Abstract-level}} \\

     & & & \multicolumn{3}{c}{\textbf{Selection-Only}} & \multicolumn{3}{c}{\textbf{Selection+Label}} & & \multicolumn{3}{c}{\textbf{Label-Only}} & \multicolumn{3}{c}{\textbf{Label+Rationale}}  \\

    \cmidrule(lr){4-6} \cmidrule(lr){7-9} \cmidrule(lr){11-13} \cmidrule(lr){14-16}

    Retrieval & Model & Row & P & R & F1 & P & R & F1 & & P & R & F1 & P & R & F1 \\

    \midrule

    \multirow{4}{*}{\shortstack[c]{\textbf{Oracle}\\ \textbf{abstract}}} &


    Oracle rationale & 1 & 100.0 & 80.5 & 89.2 & 89.6 & 72.2 & 79.9 &  & 90.1 & 77.5 & 83.3 & 90.1 & 77.5 & 83.3 \\

    \arrayrulecolor{black!30}\cmidrule(lr){2-16}

 & Zero-shot & 2 & 42.5 & 45.1 & 43.8 & 36.1 & 38.4 & 37.2 &  & 86.9 & 53.6 & 66.3 & 67.9 & 41.9 & 51.8 \\

 & \sysname & 3 & 76.1 & 63.8 & 69.4 & 66.5 & 55.7 & \textbf{60.6}          &  & 87.3 & 65.3 & 74.7 & 84.9 & 63.5 & \textbf{72.7}          \\

\thickmid

\multirow{4}{*}{\textbf{Open}} &

Oracle rationale & 4 & 100.0 & 56.5 & 72.2 & 87.6 & 49.5 & 63.2 &  & 88.9 & 54.1 & 67.2 & 88.9 & 54.1 & 67.2 \\

\arrayrulecolor{black!30}\cmidrule(lr){2-16}
 & Zero-shot & 5 & 28.7 & 37.6 & 32.5 & 23.7 & 31.1 & 26.9 &  & 56.0 & 42.3 & 48.2 & 42.3 & 32.0 & 36.4 \\

 & \sysname & 6 & 45.0 & 47.3 & 46.1 & 38.6 & 40.5 & \textbf{39.5}          &  & 47.5 & 47.3 & 47.4 & 46.6 & 46.4 & \textbf{46.5}          \\


\arrayrulecolor{black}\bottomrule

  \end{tabularx}

  \caption{Test set performance on \ours, according to the metrics from \S\ref{sec:evaluation}. For the \emph{Oracle abstract} rows, the system is provided with gold evidence abstracts. ``Oracle rationale'' rows indicate that the gold rationales are provided as input. ``Zero-shot'' indicates zero-shot performance of a verification system trained on \fever. See Appendix \ref{sec:uncertainty} for standard deviations on all estimates, and dev set performance. }
    \label{tbl:main_results}

\end{table*}


The \componenttwo module is evaluated on its ability to select rationale sentences given gold abstracts\footnote{Our \fever-trained \componenttwo module achieves 79.9 sentence-level F1 on the \fever test set, virtually identical to 79.6 reported in \citet{DeYoung2019ERASERAB}.}. The \componentthree module is evaluated on its 3-way label classification accuracy given gold rationales from cited abstracts. Cited abstracts labeled \notenough are included in the evaluation. These abstracts have no gold rationale sentences; as in \S\ref{sec:baselines}, we provide the $k$ most similar sentences from the abstract as input.

\vspace{.2cm} \noindent{\bf Results} The results are shown in Table \ref{tbl:components}. For \componentthree, the best performance is achieved by training first on the large \fever dataset and then fine-tuning on the smaller in-domain \ours training set. To understand the benefits of \fever pretraining, we examined the claim / evidence pairs where the \fever + \ours - trained model made correct predictions but the \ours - trained model did not. In  36 / 44 of these cases, the \ours - trained model predicts \notenough. Thus pretraining on \fever appears to improve the model's ability to recognize textual entailment relationships between evidence and claim -- particularly relationships indicated by non-domain-specific cues like ``is associated with'' or ``has an important role in''.

For \componenttwo, training on \ours alone produces the best results. We examined the rationales that the \ours - trained model identified but the \fever - trained model missed, and found that they generally contain science-specific vocabulary. Thus, training on additional out-of-domain data provides little benefit.

\robertalarge exhibits the strongest performance on label prediction, while \scibert has a slight edge on rationale selection. The ``claim-only'' and `'abstract-only'' models perform substantially worse than the model with access to the full input, providing reassurance that the model is not making its predictions based on statistical artifacts.



\subsection{Full task} \label{sec:main_results}
\noindent {\bf Experimental setup}
Based on the results from \S\ref{sec:pipeline_components}, we use the \componenttwo module trained on \ours only, and the \componentthree module trained on \fever + \ours for our final end-to-end system \sysname. Although \scibert performs slightly better on rationale selection, using \robertalarge for both \componenttwo and \componentthree gave the best full-pipeline performance on the dev set, so we use \robertalarge for both components. For the \componentone module, the best dev set full-pipeline performance was achieved by retrieving the top $k=3$ documents. %; adding more documents because $k > 3$ reduces end-to-end performance by introducing too many ``false positive'' candidate abstracts. See

\vspace{.2cm}\noindent{\bf Model comparisons}
We report performance of three model variants. For the ``Oracle rationale'' setting, the \componenttwo module is replaced by an oracle which outputs gold rationales for correctly retrieved documents, and no rationales for incorrect retrievals. The ``Zero-shot'' setting reports the zero-shot generalization performance of a model trained on \fever (the results on \snopes were slightly worse). \sysname reports the performance of our best system.


\vspace{.2cm}\noindent{\bf Results}
The results are shown in Table \ref{tbl:main_results}. In the oracle abstract setting, the abstract-level F1 scores are roughly comparable to label classification accuracies, %. Thus the \abstractLabelOnly F1 score in Row 1 is close to the best label prediction accuracy reported in Table \ref{tbl:components},
and the \abstractLabelRationale score in Row 3 implies an end-to-end classification accuracy of roughly 70\%, given gold abstracts.

Access to in-domain data during training clearly improves performance. Despite the small size of \ours, training on these data leads to relative improvements of 47\% on open \sentenceSelectionLabel, and 28\% on open \abstractLabelRationale (Row 6) over \fever alone.
%
The three pipeline components make similar contributions to the overall model error. Replacing \componenttwo with an oracle leads to a roughly 20-point rise in \sentenceSelectionLabel F1 (Row 6 vs. Row 4). Replacing \componentone with an oracle as well leads to a gain of roughly 20 more points (Row 4 vs. Row 1).

Nearly all correctly-labeled abstracts are supported by at least one rationale. There is only a two-point difference in F1 between \abstractLabelOnly and \abstractLabelRationale in the oracle setting (Row 3), and a one-point difference in the open setting (Row 6). The differences between \sentenceSelectionOnly and \sentenceSelectionLabel are larger, caused by examples where the model finds the evidence but fails to predict its relationship to the claim. We examine these in \S\ref{sec:error_analysis}.

%%%%%%%%%%%%%%%%%%%%%%%%%%%%%%%%%%%%%%%%%%%%%%%%%%%%%%%%%%%%%%%%%%%%%%%%%%%%%%%%

\begin{table*}[t]
  \setlength{\tabcolsep}{2pt}
  \footnotesize
  \centering
  \begin{tabular}{L{0.12\linewidth} L{0.33\linewidth} L{0.44\linewidth} R{0.06\linewidth} }
    \toprule
    \textbf{Mistake type} &  \textbf{Claim} &  \textbf{Gold evidence} & \textbf{Count} \\
    \thinmid

    Science background & Rapamycin accelerates aging in fruit flies. & \refutesCovid{\dots feeding rapamycin to adult Drosophila produces life span extension \dots} & 29 / 37  \\
    \hdashline[0.4pt/2pt]

    Directionality & Inhibiting glucose-6-phospate dehydrogenase impairs lipogenesis & \supportsCovid{\dots suppression of 6PGD decreased lipogenesis \dots} & 9 / 37 \\

    \hdashline[0.4pt/2pt]

    Numerical reasoning & Bariatric surgery reduces resolution of diabetes. & \refutesCovid{Strong associations were found \dots with a HR of 9.29 (95\% CI 6.84-12.62)...} & 9 / 37 \\

    \hdashline[0.4pt/2pt]

    Cause and effect & Major vault protein (MVP) functions to decrease tumor aggression. & \refutesCovid{Knockout of MVP leads to miR-193a accumulation...inhibiting tumor progression} & 7 / 37  \\

    \hdashline[0.4pt/2pt]

    Coreference & Low saturated fat diets have adverse effects on the development of infants & \refutesCovid{Neurological development of children in the intervention group was at least as good as ... the control group} & 6 / 37 \\

    \arrayrulecolor{black}
    \bottomrule
  \end{tabular}
  \caption{Types of reasoning required to correct 37 mistaken predictions.}
  \label{tbl:error_analysis}

\end{table*}

\subsection{Verifying claims about COVID-19} \label{sec:covid}
We conduct exploratory experiments using our system to verify claims concerning COVID-19. We tasked a medical student to write 36 COVID-related claims. For each claim $c$, we used \sysname to predict evidence abstracts $\Ehat(c)$. The annotator examined each $(c, \Ehat(c))$ pair. A pair was labeled \emph{plausible} if $\Ehat(c)$ was nonempty, and at least half of the evidence abstracts in $\Ehat(c)$ were judged to have reasonable rationales and labels. For 23 / 36 claims, the response of \sysname was deemed plausible by our annotator, demonstrating that \sysname is able to successfully retrieve and classify evidence in many cases. Two examples are shown in Table \ref{tbl:covid_example}.

% An examination of errors reveals that the system can be confused by context, where abstracts are labeled \supports or \refutes even though the rationale sentences reference a different disease or drug from the claim. An example of this is also provided in Table \ref{tbl:covid_example}.


%%%%%%%%%%%%%%%%%%%%%%%%%%%%%%%%%%%%%%%%%%%%%%%%%%%%%%%%%%%%%%%%%%%%%%%%%%%%%%%%

\subsection{Error analysis} \label{sec:error_analysis}
To better understand the errors made by \sysname, we conduct a manual analysis of 37 predictions where an evidence abstract was correctly retrieved, but where the model failed to identify any relevant rationales or predicted an incorrect label. We identify five modeling capabilities, and count the number of missed examples where each modeling capability is required to make a correct prediction (Table \ref{tbl:error_analysis}). Many examples require more than one of these capabilities.

\begin{description}[leftmargin=*,leftmargin=0pt]
  \item \textbf{Science background} includes knowledge of domain-specific lexical relationships -- e.g. Drosophila is a synonym for fruit fly.
  \item \textbf{Directionality} requires understanding increases or decreases in scientific quantities -- e.g. decreased (but not increased) lipogenesis implies impairment of the lipogenesis process.
  \item \textbf{Numerical reasoning} involves interpreting numerical or statistical statements, often reporting confidence intervals or p-values.
  \item \textbf{Cause and effect} requires reasoning about counterfactuals -- e.g. if knocking out MVP inhibits tumor progression, then the presence of MVP enables tumor growth and aggression.
  \item \textbf{Coreference} involves drawing conclusions using context stated outside of a rationale sentence -- e.g. the ``intervention group'' mentioned in Table \ref{tbl:error_analysis} was on a low saturated fat diet.
\end{description}


% TODO uncomment for camera-ready.

% When given evidence requiring an understanding of directionality or cause and effect, \sysname tended to recognize that the evidence was related to the claim, but to mix up the relationship, predicting \supports when the gold label was \refutes or vice versa. On the other hand, the model tended to not recognize evidence requiring numerical reasoning or coreference as relevant to the claim at all, and made mistakes by predicting \notenough.
